\section{Discussion and Summary}
\label{sec:discussion}

Perhaps the most useful aspect of considering system extension is that it elegantly unifies both parameter perturbations and application of harmonic forcings into a single framework. When viewed through the lens of an extended system a parameter perturbation is nothing but a harmonically forced system, albeit with a zero angular frequency (and nonlinearities). It easily extends to cases where the system parameters themselves have an oscillatory variation, which often occurs in engineering applications, with oscillating boundary conditions being a prime example. In the previous sections parameter perturbations and harmonic forcings have been treated as separate cases however, this is primarily to illustrate the identical algorithmic treatment required for the extension of the tangent space in the two cases. For the extension of the center subspace the only thing that is important is whether the extended eigenvalues are in resonance with the intrinsic critical eigenvalues of the system. For the non-resonant cases the subspace extends in a straightforward manner, while in the case of resonance one obtains a Jordan form like structure with a generalized eigenvector.

Secondly, the system extension also illuminates the structure of the problem under forcing or parameter perturbation, with the extended directions effectively being the resolvent responses of the system. Of prime importance is the evaluation of the adjoint extended modes of the system that are essential in constructing the projection operators for the individual subspace directions. The extended adjoint modes have a special structure regardless of resonance, which was fruitfully utilized to show that, when constructing asymptotic center manifold approximations of the extended system, all evaluations of the asymptotic vectors, $\AsECenM_{i_{1},\ldots,i_{k}}$, in fact only require the original linear system $\Lu$. This is of significant importantance when considering approximations for large systems where specialized algorithms are often tailored to exploit the inherent structure of $\Lu$. Despite the system extension, only the original linear system $\Lu$ is required for the calculations, which means that the specialized algorithms do not require any substantial modifications. In that sense the current approach provides the best of both worlds. It allows the calculation of center manifold approximations of perturbed systems using the original unperturbed linear operators. At the same time it provides the theoretical perspective of understanding the perturbations as intrinsic modes of an extended system. 

The upper triangular structure of the $\Khat$ matrix, which is the linear part of the reduced representation $\Param(\vecparv)$, is responsible for the cross coupling of vectors $\AsCenM_{i_{i},\ldots,i_{k}}$ at each order $k$ in the asymptotic approximations. However, at each order, the system of equations generated for the unknown $\AsCenM_{i_{i},\ldots,i_{k}}$ is lower triangular and can be solved one vector at a time through back substition. This lower triangular nature of the system of equations is also recognized by \citet{vizzaccaro24}, where the authors note that through a suitable ordering the system can be solved iteratively. By considering the ordering of the index set as defined in equation~\eqref{eqn:index_set}, this ordering is directly obtained. For large systems this is an important property since the coupled system solution can quickly become computationally infeasible.

 The derived center manifold approximation is applied to a model Ginzburg-Landau problem at the bifurcation point. The system is considered with perturbations of the diffusion coefficient $(\delta_{4})$ as well as with resonant and non-resonant external forcings. In all cases a very good agreeement is found between the nonlinear system response and the response predicted by the calculated center manifold approximations. In addition, the asymptotic center manifold is also evaluated via the graph method presented earlier, which allows for an effecient and straightforward mapping from an instantaneous nonlinear solution $\evecu(t)$ to the parameterized coordinates $\vecparv(t)$ ``on the fly''. This simple reconstruction via projection is only possible if one has knowledge of the adjoint vectors $\EAevMatrix$ in the extended space. 
 
 Finally, one may note that in the current study, only linear forcing was considered, so that $\fieldf^{H}(\vectheta)$ defined in equation~\eqref{eqn:harmonic_forcing} only has terms with linear dependence on $\vectheta$. However, this term may easily be modified to include nonlinear terms of the type $\vectheta\vectheta$, $\vectheta\vecu$, \etc This makes no difference in the overall procedure except that the nonlinear terms $\NewNonLin(\vecu,\vecnu,\vectheta)$ need to include the added nonlinearities. The rest of the procedure for tangent space extension and center manifold evaluation follows exactly as before. In fact, since we now understand that $\vecnu$ and $\vectheta$ can be treated equivalently, the parameter perturbation term $\delta'_{4}\nabla^{2}A$ is precisely one such term. If one endows $\delta'_{4}$ with an oscillatory variation one obtains what \citet{thurnher24} refer to as parametric forcing. From the point of the view of the extended system $\mathfrak{Im}(\lambda^{P}_{i})$ being non-zero makes no algorithmic difference. 

\backmatter

\bmhead{Acknowledgements}
The research was supported by the Okinawa Institute of Science and Technology Graduate	University (OIST) with subsidy funding from the Cabinet Office, Government of Japan. 












